\scanpage{23}

\uline{Théorème 2} Soit $f:X\longrightarrow Y$ un morphisme \deleted{loc. de \unsure{t}} propre de préschémas \inserted{\unsure{loc}} \unsure{noeths|noeth}, avec $Y$ intègre, $y_0\in Y$, $y_1$ \deleted{le pt gén. de $Y$. \inserted{\unsure{une générisation}}} \inserted{le pt gén de $Y$}. On suppose \deleted{que l'on est dans le \unsure{l'un des cas suivants}} \deleted{\unsure{intég.}} \deleted{\unsure{$y_1$ est un pt gén}}%
% vertical rule beside the first three lines of the hypotheses

\margin{?}\deleted{(\unsure{a})} (a) $Y$ est [\unsure{géom.}\textsuperscript{t}] unibranche en $y_0$, \deleted{\unsure{et}} tte composante irréductible de $X$ domine $Y$.

\deleted{(b) $f$ est universellt ouvert} \inserted{(b) les composantes irréductibles de} \deleted{\unsure{ses en}} $X_1$ sont de dim $\geq k+1$, et $X_1$ est connexe en dim $\geq k$.

Alors (les composantes irréductibles de $X_0$ sont de dim $\geq k+1$) et $X_0$ est \uline{connexe} en dim $\geq k$.

L'assertion entre parenthèses est bien connue par le théorème des 2 dimensions, et n'utilise pas l'hypothèse unibranche. Celle-ci sert uniquemt à appliquer le \uline{th. de connexion}, qui assure que $X_0$ est connexe.

% archive sheet number: 22
